% ==================================================================
% id:       comb-015
% topic:    קומבינטוריקה
% subtopic: עקרון הכפל וסידורים
% solution: full
% ==================================================================
כמה מספרים בני 5 ספרות אפשר לרשום מהספרות $1,2,3,4,5,7,9$, כך ש:
\begin{enumerate}
\item המספר מתחלק ב-5, והספרות לא חוזרות על עצמן?
\begin{solbox}
מספר מתחלק ב-5 אם ספרת האחדות שלו היא $0$ או $5$. בין הספרות שלנו אין $0$, ולכן ספרת האחדות היא $5$. את ארבע הספרות האחרות בוחרים מתוך 6 הספרות שנשארו, עם משמעות לסדר ובלי חזרות: $6\cdot5\cdot4\cdot3=\frac{6!}{2!}=\ans{360}$.
\end{solbox}
\item המספר מתחלק ב-5, ומותרת חזרה על הספרות?
\begin{solbox}
שוב ספרת האחדות היא $5$, ובכל אחד מארבעת המקומות האחרים יש 7 אפשרויות: $7^4=\ans{2401}$.
\end{solbox}
\end{enumerate}
